\section{Bounded division and an auxiliary Noetherian ring}
\label{sec:division}

Put
\begin{equation}\label{eq:B}
 F=k[[x]],\quad L=k((x)),\quad A_0=F[U],\quad
 S=k[U]\setminus\{0\},\quad B=S^{-1}A_0.
\end{equation}
The natural embeddings identify $B$ with $K\otimes_kF$ inside
$K[[x]]$. The ring $B$ is a Noetherian domain, and
\begin{equation}\label{eq:Bdim}
 \dim B\le r.
\end{equation}
Indeed, if $\q\in\Spec B$ avoids $x$, its lower prime chains lie
in $B[1/x]$, a localization of $L[U]$, and hence $\height\q\le r$.
If $x\in\q$, then $B/xB=K$ gives $\q=xB$. The principal ideal
theorem gives $\height(xB)=1\le r$.

Extend the filtration to $B$ by
\begin{equation}\label{eq:E}
 E_N=W_N\otimes_kF
 =\left\{\frac{p}{D_N}:p\in F[U],\
                    \deg p\le\deg D_N+N\right\}.
\end{equation}
Under the embedding into $K[[x]]$, one also has
\begin{equation}\label{eq:finite-basis}
 W_N[[x]]=E_N.
\end{equation}
Here the left side means series whose coefficients lie in $W_N$.
To verify the identity, expand each coefficient in the finite
$k$-basis $\{U^\alpha/D_N:|\alpha|\le\deg D_N+N\}$ and collect
the coefficients of each basis element into a series in $F$.
The $E_N$ exhaust $B$, are increasing and multiplicative, and define
$w(b)=\min\{N:b\in E_N\}$ on $B$. The inequalities
\eqref{eq:weight} still hold, and every element of $F$ has weight zero.

\begin{lemma}[Bounded division]\label{lem:division}
Fix $b_1,\ldots,b_s\in B$ and let $J=(b_1,\ldots,b_s)B$.
There is an integer $C\ge0$ such that, for every $N\ge0$ and
$h\in J\cap E_N$, one can write
\begin{equation}\label{eq:division}
 h=\sum_{i=1}^{s}b_i z_i,\qquad z_i\in E_{N+C}.
\end{equation}
\end{lemma}

\begin{proof}
The case $J=0$ is immediate. Set $J_0=J\cap A_0$ and filter
$A_0=F[U]$ by total degree in $U$. Its associated graded ring is
$F[U]$, which is Noetherian. Thus the homogeneous ideal generated
by the highest-degree parts of elements of $J_0$ is generated by
the highest-degree parts of finitely many nonzero elements
$\ell_1,\ldots,\ell_t\in J_0$. Write $d_j=\deg\ell_j$.
Every $0\ne a\in J_0$ admits a representation
\begin{equation}\label{eq:degree-division}
 a=\sum_{j=1}^{t}\ell_jp_j,\qquad
 p_j\in F[U],\qquad \deg p_j\le\deg a-d_j.
\end{equation}
Indeed, express the highest-degree part of $a$ using homogeneous
multipliers of degree $\deg a-d_j$, subtract the corresponding
combination of the $\ell_j$, and repeat at the smaller degree.
This process terminates. A negative degree bound means $p_j=0$;
no division by a nonunit of $F$ is involved.

Since $\ell_j\in J$, there are fixed $s_0\in S$ and
$A_{ij}\in F[U]$ with
\[
 \ell_j=\sum_{i=1}^{s}b_i\frac{A_{ij}}{s_0}
 \qquad(1\le j\le t).
\]
For $h\in J\cap E_N$, write $h=a/D_N$, where
$a\in A_0$ and $\deg a\le\deg D_N+N$.
Since $D_N$ is a unit of $B$, we have $a\in J_0$.
Using~\eqref{eq:degree-division}, set
\[
 z_i=\frac{\sum_j A_{ij}p_j}{s_0D_N};
 \qquad h=\sum_i b_i z_i.
\]
There is a fixed $C_0\ge0$ such that each numerator here has degree
at most $\deg D_N+N+C_0$.

Choose $C$ so large that
\begin{equation}\label{eq:denominator-shift}
 s_0D_N\mid D_{N+C}\quad(N\ge0),\qquad C\ge C_0-\deg s_0.
\end{equation}
To obtain the divisibility, for each factor $q_j^{e_j}$ of $s_0$
it suffices to require $C\ge je_j$, since
\[
 \left\lfloor\frac{N+C}{j}\right\rfloor
 -\left\lfloor\frac{N}{j}\right\rfloor\ge e_j.
\]
Writing $z_i$ over $D_{N+C}$ now bounds its numerator degree by
\[
 \deg D_N+N+C_0+\deg\frac{D_{N+C}}{s_0D_N}
 =\deg D_{N+C}+N+C_0-\deg s_0
 \le\deg D_{N+C}+N+C.
\]
Thus $z_i\in E_{N+C}$, as required.
\end{proof}

Inside $B[[c]]$, define
\begin{equation}\label{eq:H}
 H=\left\{\sum_{n\ge0}b_nc^n:
       \text{there is }M\ge0\text{ with }w(b_n)\le M2^n
       \text{ for all }n\right\}.
\end{equation}
The bound may depend on the series. This is a subring containing
$B$: for multiplication, the weight of the coefficient of $c^n$
is at most
\[
 \max_{0\le i\le n}(M2^i+M'2^{n-i})\le(M+M')2^n.
\]
Shifting the coefficients gives
\begin{equation}\label{eq:Hgraded}
 H\cap c^eB[[c]]=c^eH\quad(e\ge0),\qquad
 H/cH\cong B,\qquad \gr_cH\cong B[z],
\end{equation}
where $z$ is the initial form of $c$. Furthermore,
\begin{equation}\label{eq:jacobson}
 c\in\Jac(H).
\end{equation}
Indeed, for every $h\in H$, the positive-degree coefficients of
$ch$ satisfy an exponential bound. Formula~\eqref{eq:reciprocal}
shows that $(1-ch)^{-1}\in H$.

\Needspace{6\baselineskip}
\begin{proposition}\label{prop:noetherian}
The ring $H$ is Noetherian, and $\dim H[1/c]\le r$.
\end{proposition}

\begin{proof}
Let $0\ne I\subseteq H$ be an ideal. Since $B[z]$ is Noetherian,
the ideal of initial forms of $I$ is generated by the initial forms
of finitely many nonzero elements $f_1,\ldots,f_s\in I$. Write
\[
 e_i=\ord_c(f_i),\qquad
 f_i=\sum_{j\ge0}f_{i,e_i+j}c^{e_i+j},\qquad b_i=f_{i,e_i}.
\]
For each $n\ge0$ set $J_n=(b_i:e_i\le n)B$.
There are only finitely many such generator lists, so
\cref{lem:division} supplies one constant $C\ge0$ valid for
all the nonzero $J_n$.

Fix $g=\sum_{n\ge0}g_nc^n\in I$. Choose $A,F_0\ge0$ with
\begin{equation}\label{eq:input-bounds}
 w(g_n)\le A2^n\quad(n\ge0),\qquad
 w(f_{i,e_i+j})\le F_0 2^j\quad(j\ge1,\ 1\le i\le s).
\end{equation}
We cancel the coefficients of $g$ successively. After canceling
the coefficients below degree $n$ by finitely many multiples of
the $f_i$, the residual still lies in $I$. Its coefficient $h_n$
at degree $n$ belongs to $J_n$: if nonzero, its initial form lies
in the homogeneous ideal generated by $b_i z^{e_i}$.
By bounded division choose $a_{i,n}\in B$ such that
\begin{equation}\label{eq:cancellation}
 h_n=\sum_{e_i\le n}b_i a_{i,n},\qquad
 w(a_{i,n})\le w(h_n)+C,
\end{equation}
and put $a_{i,n}=0$ when $e_i>n$.
If $h_n=0$, choose all coefficients to be zero.
Subtracting $\sum_{e_i\le n}a_{i,n}c^{n-e_i}f_i$ cancels
degree $n$ as well.

We must verify that the resulting infinite multipliers belong to
$H$. Put $T_n=\max_i w(a_{i,n})$. A cancellation at degree
$m<n$ contributes to degree $n$ a product
$a_{i,m}f_{i,e_i+n-m}$. Thus~\eqref{eq:weight},
\eqref{eq:input-bounds}, and~\eqref{eq:cancellation} give
\begin{equation}\label{eq:recurrence}
 T_n\le\max\left\{A2^n,
       \max_{0\le m<n}\bigl(T_m+F_0 2^{n-m}\bigr)\right\}+C,
\end{equation}
where the inner maximum is omitted for $n=0$.
Choose $M\ge\max\{A+C,\,2F_0+C\}$.
We claim $T_n\le M2^n$. The case $n=0$ is immediate.
If the claim holds below $n$, then, for $d=n-m\ge1$,
\begin{align*}
 M(2^n-2^m)
 &=M2^m(2^d-1)\\
 &\ge(2F_0+C)(2^d-1)\ge F_0 2^d+C.
\end{align*}
It follows that $T_m+F_0 2^{n-m}+C\le M2^n$; also
$A2^n+C\le M2^n$. Equation~\eqref{eq:recurrence} proves
the induction step. Consequently
\[
 h_i=\sum_{n\ge e_i}a_{i,n}c^{n-e_i}\in H,
\]
because its coefficient of degree $j$ has weight at most
$M2^{e_i}2^j$. The cancellation gives the coefficientwise identity
$g=\sum_i f_i h_i$ in $B[[c]]$, and hence in $H$.
Thus $I=(f_1,\ldots,f_s)H$. The zero ideal is finitely generated
as well, proving that $H$ is Noetherian.

For every maximal ideal $\mathfrak n$ of $H$,
\eqref{eq:jacobson} and the Noetherian local dimension inequality
for a principal quotient give
\[
 \dim H_{\mathfrak n}
 \le\dim(H_{\mathfrak n}/cH_{\mathfrak n})+1
 \le\dim B+1\le r+1.
\]
Therefore $\dim H\le r+1$. Every prime chain avoiding $c$ can
be extended strictly by a maximal ideal, since every maximal ideal
contains $c$. Thus every such chain has length at most $r$, and
\mbox{$\dim H[1/c]\le r$}.
\end{proof}
